\documentclass[12pt,a4paper]{article}
\usepackage[UTF8]{ctex}
\usepackage{geometry}
\usepackage{amsmath,amsfonts,amssymb}
\usepackage{booktabs,longtable}
\usepackage{tabularx}
\usepackage{adjustbox}
\usepackage{graphicx}
\usepackage{float}
\usepackage{hyperref}
\usepackage{bookmark}
\usepackage{caption}
\usepackage{enumitem}
\usepackage{listings}
\usepackage{xcolor}
\usepackage[numbers]{natbib}

\geometry{left=2.5cm,right=2.5cm,top=2.5cm,bottom=2.5cm}

\setmonofont{Menlo}[Scale=MatchLowercase]

% ---------- 代码 listings 设置 ----------
\lstset{
    language=Python,
    basicstyle=\ttfamily\small,
    keywordstyle=\color{blue},
    commentstyle=\color{gray},
    stringstyle=\color{orange},
    numbers=left,
    numberstyle=\tiny\color{gray},
    frame=single,
    breaklines=true,
    tabsize=4
}

% ---------- 参考文献格式：GB/T 7714-2015 ----------
\bibliographystyle{gbt7714-numerical}

\title{\textbf{中老年人群高血脂症的风险预警及干预方案优化}}
\author{}
\date{}

\begin{document}

% 首页（题目、摘要、关键词）不显示页码
\pagestyle{empty}
\maketitle

\begin{abstract}
本文针对中老年人群高血脂症的风险预警及干预方案优化问题，融合中医体质学、血常规体检指标及日常活动能力评分，构建了多维度数学模型。\\[6pt]
\textbf{针对问题1}，构建``Pearson/Spearman相关+组间差异检验+逐步Logistic回归''的三层筛选体系，从血常规与活动量表中筛选出TC、TG、HDL-C、尿酸、LDL-C为预警高血脂的关键指标；同时发现血常规指标与痰湿质积分的线性关联极弱（调整$R^2\approx0$），提示痰湿体质的判定更依赖中医量表而非单一生化指标。通过多元Logistic回归量化九种体质对发病风险的贡献度，气虚质与平和质的系数趋势相对较大，但均未达到显著性水平，可能与体质积分间的耦合效应及样本量限制有关。\\[6pt]
\textbf{针对问题2}，基于Logistic回归概率输出构造综合风险评分，采用``连续数据离散化''思想结合临床锚点与统计分位数确定低/中/高三级风险阈值；利用5折分层交叉验证与ROC曲线量化模型区分能力（AUC=0.777），并进行权重敏感性分析；通过条件概率识别核心特征组合，如``痰湿积分$\geq 60$ \& 活动能力$<40$''等。\\[6pt]
\textbf{针对问题3}，建立带依从性约束的组合优化模型，以中医调理原则强制匹配调理分级，在年龄、活动能力耐受度及老年干预疲劳效应约束下，遍历活动强度与训练频率的所有合法组合，最小化6个月后痰湿积分。求解得到样本ID 1、2、3的最优干预方案，并总结出``患者特征—最优方案''匹配规律。\\[6pt]
\textbf{关键词：} 高血脂症；中医体质；Logistic回归；风险分层；整数规划；干预优化
\end{abstract}


\newpage
% 目录页不显示页码
\tableofcontents

\newpage

% ---------- 正文开始，页码从 1 计 ----------
\pagestyle{plain}
\setcounter{page}{1}

% ==========================================================
\section{问题重述与分析}
% ==========================================================

\subsection{问题背景}
人口老龄化加剧，中老年人群高血脂症发病率攀升，成为心脑血管疾病的重要危险因素。中医体质学认为痰湿体质与血脂代谢异常高度契合，而日常活动能力直接影响气血运行与膏脂运化~\cite{wang2009tcm}。目前临床筛查单一依赖血常规，缺乏对体质分型与活动能力的综合考量，难以实现``治未病''的精准防控。

\subsection{问题提出}
\begin{enumerate}[label=\textbf{问题\arabic*}]
    \item 从血常规与活动量表中筛选能有效表征痰湿体质严重程度、且预警高血脂发病风险的关键指标；研究九种体质对发病风险的贡献度差异。
    \item 构建融合多维度特征的三级风险预警模型（低/中/高），明确阈值选取依据；识别痰湿体质高风险人群的核心特征组合。
    \item 针对痰湿体质患者，在中医调理原则、身体耐受度及经济成本约束下，构建6个月干预方案优化模型，给出样本ID 1、2、3的最优方案及``患者特征—最优方案''匹配规律。
\end{enumerate}

\subsection{问题分析}
\begin{itemize}
    \item \textbf{问题1}属于特征筛选与归因分析问题。可利用相关性分析与回归模型的显著性检验筛选关键指标；利用Logistic回归系数或优势比(OR)衡量体质贡献度。
    \item \textbf{问题2}属于风险分层与规则挖掘问题。核心是将连续风险概率离散化为三级，需兼顾统计分布与临床意义；核心特征组合可通过条件概率或决策树规则提取。
    \item \textbf{问题3}属于带约束的组合优化问题。决策空间极小（调理分级$\times$活动强度$\times$频率），可用穷举法精确求解；目标是降分最大化与成本可控的双重要求。
\end{itemize}

% ==========================================================
\section{模型假设与符号说明}
% ==========================================================

\subsection{基本假设}
\begin{enumerate}
    \item 样本数据独立同分布，能代表目标中老年人群总体特征。
    \item 中医调理分级严格按痰湿积分区间匹配，不可跨级选择。
    \item 活动干预效果满足题目给出的线性叠加规则：基础强度效果与额外频率效果相互独立。
    \item 每月按4周计，6个月干预周期共24周。
    \item 患者对干预方案的依从性为$100\%$，不考虑中途退出。
    \item 高龄（80--89岁）及低活动能力（$M<40$）患者存在干预疲劳效应：每周超过5次的部分，实际有效频率依从性衰减，且设置安全上限$f_{\max}=7$次/周。
\end{enumerate}

\subsection{符号说明}

\begin{table}[H]
\centering
\caption{主要符号及其含义}
\begin{tabular}{clc}
\toprule
\textbf{符号} & \textbf{含义} & \textbf{单位/备注} \\
\midrule
$x$ & 中医调理分级 & $x\in\{1,2,3\}$ \\
$y$ & 活动干预强度 & $y\in\{1,2,3\}$ \\
$f$ & 每周训练频率 & $f\in\{1,2,\dots,10\}$（次/周） \\
$S_0$ & 初始痰湿积分 & $0\sim100$分 \\
$S_6$ & 6个月后痰湿积分 & $0\sim100$分 \\
$r$ & 每月痰湿积分下降率 & $0\sim$某正值 \\
$C_T(x)$ & 调理分级$x$的月成本 & $\{30,80,130\}$元/月 \\
$C_A(y)$ & 活动强度$y$的单次成本 & $\{3,5,8\}$元/次 \\
$M$ & 活动量表总分(ADL+IADL) & $0\sim100$分 \\
$A$ & 年龄组 & $1\sim5$（40$\sim$89岁） \\
$P$ & 高血脂发病概率 & $0\sim1$ \\
$R$ & 综合风险评分 & 用于三级分层 \\
\bottomrule
\end{tabular}
\end{table}

% ==========================================================
\section{问题1：关键指标筛选与体质贡献度}
% ==========================================================

\subsection{数据预处理}
原始数据共1000例样本，包含9种体质积分、血常规指标（HDL-C、LDL-C、TG、TC、空腹血糖、血尿酸）、BMI、ADL/IADL量表及年龄、性别等人口学特征。为保证建模质量，预处理分为以下三步：

\textbf{（1）数据整合与变量映射}
原始数据为CSV格式，共1000例样本、36个字段。首先对列名进行标准化映射，将中文字段名转为英文简写以便后续建模。依据题目附表1的编码规则，将性别（0=女，1=男）、年龄组（1=40--49岁至5=80--89岁）等分类变量转为数值型，确保类型一致。

\textbf{（2）血脂异常项数构造}
依据题目给定的临床正常范围，逐样本统计TC、TG、LDL-C、HDL-C四项血脂指标中超出范围的项数，记为$N_{abn}\in\{0,1,2,3,4\}$。血尿酸正常范围存在性别差异（男$208\sim428\,\mu\text{mol/L}$，女$155\sim357\,\mu\text{mol/L}$），因此单独构造血尿酸异常标记$U_{abn}\in\{0,1\}$，不混入血脂异常项数，以保证指标定义的临床准确性。

\textbf{（2）Z-score标准化}
由于血常规指标量纲各异（如TC单位为mmol/L，BMI为kg/m$^2$），直接入模会导致大数量级变量压制小数量级变量。因此，对TC、TG、LDL-C、HDL-C、血糖、尿酸、BMI、ADL、IADL及活动量表总分进行标准化：
\begin{equation}
Z_{ij}=\frac{X_{ij}-\bar{X}_j}{S_j}
\end{equation}
其中$\bar{X}_j$与$S_j$分别为第$j$个变量的样本均值与标准差。

\textbf{（3）血脂异常项数构造}
依据《中国成人血脂异常防治指南（2016年修订版）》~\cite{zhinan2016}，定义血脂正常范围为：TC $3.1\sim6.2$ mmol/L，TG $0.56\sim1.7$ mmol/L，LDL-C $2.07\sim3.1$ mmol/L，HDL-C $1.04\sim1.55$ mmol/L。逐样本统计超出上述范围的项数，记为$N_{abn}\in\{0,1,2,3,4\}$，作为后续风险评分的合成指标之一。

\subsection{关键指标筛选方法}

\subsubsection{相关性分析与组间差异检验}
题目要求筛选的指标需同时满足两个功能：\textbf{表征痰湿体质严重程度}与\textbf{预警高血脂发病风险}。因此，本文从两个目标变量出发，建立三层分析体系：

\textbf{（1）Pearson与Spearman相关分析}
针对痰湿质积分（连续变量），分别计算Pearson积矩相关系数（捕捉线性关系）与Spearman秩相关系数（捕捉单调非线性关系）。针对高血脂二分类标签（0/1），计算点二列相关系数$r_{pb}$，其公式为：
\begin{equation}
r_{pb}=\frac{\bar{X}_1-\bar{X}_0}{S_X}\sqrt{\frac{n_1n_0}{n(n-1)}}
\end{equation}
其中$\bar{X}_1,\bar{X}_0$分别为患病组与未患病组的均值，$S_X$为总体标准差。

\textbf{（2）组间差异检验（Mann-Whitney U检验）}
由于痰湿质积分在样例数据中与血常规指标的线性相关性极弱（见后文结果），单一相关分析可能遗漏非线性关联。为此，引入基于体质标签的组间差异检验：将样本按主体质是否为痰湿质（体质标签$=5$）分为两组，对各血常规及活动量表指标进行Mann-Whitney U检验。该检验不假设正态分布，适用于检测两组间的分布差异。效应量采用秩次相关系数$r=Z/\sqrt{N}$衡量。

\textbf{（3）多元线性回归（控制混杂因素）}
进一步以痰湿质积分为因变量，10项血常规及活动量表指标为自变量，建立多元线性回归模型，通过标准化系数（Beta系数）比较各指标在控制其他变量后的独立贡献。

\subsubsection{逐步Logistic回归}
相关性分析仅能揭示变量间的线性关联，无法排除混杂因素，因此进一步建立Logistic回归模型，通过控制其他变量来识别独立预警指标。设$Y\in\{0,1\}$为高血脂标签，$\boldsymbol{X}=(X_1,X_2,\dots,X_k)^T$为候选指标向量，Logistic模型形式为：
\begin{equation}
\ln\frac{P(Y=1|\boldsymbol{X})}{1-P(Y=1|\boldsymbol{X})} = \beta_0 + \sum_{i=1}^{k}\beta_i X_i
\end{equation}
其中左侧为事件发生比（Odds）的对数，称为logit函数；$\beta_i$表示在控制其他变量后，$X_i$每增加一个单位对logit的边际贡献。

为从高维候选指标中自动筛选最优子集，采用\textbf{前向选择（Forward Selection）}与\textbf{后向剔除（Backward Elimination）}相结合的逐步策略：
\begin{enumerate}
    \item \textbf{前向选择}：从未选变量中找到$p\text{-value}<\alpha_{in}=0.05$的变量，加入模型；
    \item \textbf{后向剔除}：对已选变量进行Wald检验，若某变量$p\text{-value}>\alpha_{out}=0.10$，则将其剔除；
    \item 重复上述两步，直至没有变量可进也没有变量可出，算法收敛。
\end{enumerate}

该策略相比全子集遍历大幅降低了计算量，同时通过双向校验减少了单一方向选择可能遗漏的重要变量。最终保留的变量即为在统计学意义上能够独立预警高血脂的关键指标。

\subsection{九种体质贡献度分析}
中医体质学将人体划分为九种基本类型，不同体质与疾病易感性存在差异~\cite{wang2009tcm}。为量化各体质对高血脂发病风险的独立贡献，将9种体质积分同时纳入多元Logistic回归。由于体质积分间可能存在相关性（如痰湿质与湿热质常相伴出现），直接比较原始系数会因量纲差异而产生偏误。因此，先对各体质积分进行Z-score标准化：
\begin{equation}
X_i^* = \frac{X_i-\bar{X}_i}{\sigma_i}, \qquad i=1,2,\dots,9
\end{equation}
标准化后，各变量均值为0、标准差为1，回归系数$\beta_i^*$可直接比较绝对值大小。

模型形式为：
\begin{equation}
\text{logit}(P) = \beta_0 + \sum_{i=1}^{9}\beta_i^* X_i^*
\end{equation}

在解释贡献度时，综合考察两项指标：
\begin{itemize}
    \item \textbf{标准化回归系数}$|\beta_i^*|$：反映该体质积分每增加1个标准差，logit的变化幅度；
    \item \textbf{优势比（OR）}$\text{OR}_i=\exp(\beta_i^*)$：反映该体质积分每增加1个标准差，患病优势（Odds）的倍数变化。OR>1表示风险增加，OR<1表示风险降低。
\end{itemize}
需要注意的是，当某样本的体质标签（主体质）与特定体质积分高度相关时，可能引入多重共线性，导致个别体质系数的标准误膨胀。此时应结合单因素方差分析（ANOVA）或VIF（方差膨胀因子）诊断，综合判断贡献度排序。

\subsection{问题1结果与分析}

\subsubsection{相关性分析与组间差异检验结果}
表~\ref{tab:corr}展示了各指标与痰湿质积分及高血脂标签的相关系数。在\textbf{预警高血脂}方面，TG、TC的点二列相关系数最高（分别为0.467、0.439，$p\ll0.001$），LDL-C、HDL-C及尿酸也显示出中等强度关联，与临床认知一致。然而，在\textbf{表征痰湿体质严重程度}方面，所有血常规及活动量表指标与痰湿质积分的Pearson及Spearman相关系数绝对值均低于0.06，多元线性回归的调整$R^2$仅为$-0.004$，无一指标通过显著性检验；Mann-Whitney U检验亦未发现痰湿体质组与非痰湿体质组在各指标上存在显著分布差异（效应量$|r|<0.05$）。

这一结果具有重要的方法论启示：\textbf{痰湿体质的判定主要依赖中医体质辨识量表（如舌苔、脉象、症状问诊）而非单一血生化指标}，血常规指标更适合作为``痰湿体质是否已继发血脂异常''的下游监测工具，而非表征痰湿严重程度的直接度量。因此，本文采用``多数投票''综合筛选策略——凡在三个分析维度中至少满足一项（Spearman $|r|>0.10$、组间检验$p<0.05$且$|r|>0.05$、或高血脂点二列$|r|>0.15$）的指标均纳入关键指标集，最终入选TG、TC、LDL-C、HDL-C、尿酸五项。

\begin{table}[H]
\centering
\caption{关键指标与目标变量的相关系数}
\label{tab:corr}
\begin{tabular}{lccc}
\toprule
\textbf{指标} & \textbf{与痰湿质Pearson $r$} & \textbf{与高血脂点二列 $r$} & \textbf{综合入选} \\
\midrule
TG & $-0.020$ & \textbf{0.467}$^{***}$ & $\checkmark$ \\
TC & $-0.042$ & \textbf{0.439}$^{***}$ & $\checkmark$ \\
LDL-C & 0.017 & \textbf{0.189}$^{***}$ & $\checkmark$ \\
HDL-C & 0.015 & $-0.177^{**}$ & $\checkmark$ \\
尿酸 & $-0.003$ & 0.162$^{**}$ & $\checkmark$ \\
BMI & $-0.006$ & 0.026 & --- \\
活动量表总分 & $-0.012$ & 0.020 & --- \\
\bottomrule
\multicolumn{4}{l}{\small 注：$^{*}p<0.05$, $^{**}p<0.01$, $^{***}p<0.001$}
\end{tabular}
\end{table}

\subsubsection{逐步回归筛选结果}
逐步Logistic回归在10个候选指标中最终保留5个：TC、TG、HDL-C、尿酸、LDL-C。BMI与活动能力指标因未能通过进入阈值（$p\geq0.05$）而被排除。这表明在控制其他血常规指标后，BMI对高血脂的独立增量信息有限，其效应可能已被血脂指标所吸收。

\subsubsection{体质贡献度排序}
标准化后的多元Logistic回归结果如表~\ref{tab:const}所示。值得注意的是，\textbf{九种体质的标准化系数均未达到传统显著性水平}（$p\text{-value}\geq0.12$），提示在同时控制其他体质积分后，单一体质对高血脂的独立边际贡献有限。从系数趋势看，气虚质（$\beta^*=0.123$，$\text{OR}=1.130$）和平和质（$\beta^*=0.083$，$\text{OR}=1.087$）的绝对值相对较大，但二者的95\%置信区间均包含0，不宜做过度解读。

痰湿质 surprisingly 排名靠后（$\beta^*=-0.008$，$\text{OR}=0.992$，$p=0.919$），这一反直觉现象的可能原因在于：\textbf{体质标签（主体质）与体质积分高度耦合}，当主体质已被标签变量部分解释后，痰湿积分的边际贡献被压缩；此外，样例数据中体质积分与血脂指标的线性关联整体偏弱，也可能导致回归系数不稳定。进一步计算VIF值（表~\ref{tab:const}），九种体质的VIF均在$2.1\sim3.6$之间，远低于5的常用阈值，说明多重共线性并非导致系数不显著的主因。若后续研究能获取更大样本或补充痰湿体质专属生化标志物（如脂联素、炎症因子），有望提升模型的统计效能。

\begin{table}[H]
\centering
\caption{九种体质对高血脂发病风险的标准化Logistic回归结果}
\label{tab:const}
\begin{tabular}{lcccc}
\toprule
\textbf{体质} & \textbf{标准化系数$\beta^*$} & \textbf{OR} & \textbf{$p$-value} & \textbf{VIF} \\
\midrule
气虚质  & 0.123 & 1.130 & 0.124 & 3.31 \\
平和质  & 0.083 & 1.087 & 0.340 & 2.07 \\
湿热质  & $-0.051$ & 0.951 & 0.524 & 3.34 \\
血瘀质  & $-0.048$ & 0.954 & 0.550 & 3.59 \\
阳虚质  & 0.021 & 1.021 & 0.796 & 3.20 \\
痰湿质  & $-0.008$ & 0.992 & 0.919 & 3.39 \\
\bottomrule
\multicolumn{5}{l}{\small 注：剩余三种体质（气郁质VIF=3.47、特禀质VIF=3.29、阴虚质VIF=3.31）系数绝对值均$<0.02$，$p>0.85$，从略。}
\end{tabular}
\end{table}

% ==========================================================
\section{问题2：三级风险预警模型}
% ==========================================================

\subsection{综合风险评分构造}
单一血常规指标或单一体质积分均难以全面刻画高血脂风险。为此，本文构建融合\textbf{血常规发病概率}、\textbf{痰湿体质严重程度}与\textbf{血脂异常项数}的三维综合风险评分。

首先，以问题1筛选的5个关键指标（TC、TG、HDL-C、尿酸、LDL-C）为特征，训练Logistic回归模型，输出每个样本的基础发病概率$P_i\in[0,1]$。该概率反映了在血常规维度上的患病风险。

其次，将痰湿质积分$S_{phlegm}$（范围$0\sim100$）和血脂异常项数$N_{abn}$（范围$0\sim4$）分别归一化到$[0,1]$区间，构造综合风险评分：
\begin{equation}
R_i = w_1\cdot P_i + w_2\cdot\frac{S_{phlegm}}{100} + w_3\cdot\frac{N_{abn}}{4}
\end{equation}

权重$w_1,w_2,w_3$的确定有两种思路：
\begin{itemize}
    \item \textbf{熵权法}：根据各指标的信息熵客观赋权，熵越小（变异越大）权重越高；
    \item \textbf{专家经验法}：鉴于血常规概率$P_i$直接对应临床诊断逻辑，赋予最高权重0.5；痰湿质积分反映中医体质维度，权重0.3；血脂异常项数为辅助确认指标，权重0.2。本文采用后者，以保证评分体系与临床认知一致。
\end{itemize}

该评分将$[0,1]$区间的概率、$[0,100]$区间的体质积分、$\{0,1,2,3,4\}$的计数指标统一到$[0,1]$尺度，便于后续离散化分层。为验证权重选取的稳健性，第~\ref{sec:sens}节进行了敏感性分析。

\subsection{三级风险阈值确定}
综合风险评分$R_i$为连续变量，取值范围约为$[0.15,0.95]$。为将其转化为临床可操作的低/中/高三级预警，需要确定两个阈值$\theta_L$（低$\to$中）与$\theta_H$（中$\to$高）。本文提出\textbf{统计分位数+临床锚定}的双轨制阈值确定方法，兼顾数据分布规律与医学先验知识。
\subsubsection{统计分位数法}
分位数法依据样本自身的分布特征划分层级，不依赖于外部标准：
\begin{itemize}
    \item 高风险阈值$\theta_H$：取$R_i$的75\%分位数，即样本中风险评分最高的25\%人群被划分为高风险；
    \item 低风险阈值$\theta_L$：取$R_i$的30\%分位数，即样本中风险评分最低的30\%人群被划分为低风险；
    \item 中间40\%为中等风险。
\end{itemize}
分位数法的优点在于自适应数据分布，避免了主观设定阈值可能导致的极端倾斜；其局限在于完全数据驱动，可能将临床意义上的高危患者漏划至中风险。
\subsubsection{临床锚定法}
为弥补纯数据驱动的不足，引入题目给定的临床先验作为锚点：
\begin{enumerate}
    \item 若样本血脂指标异常（$N_{abn}\geq1$）且痰湿积分$\geq60$，则该样本在医学上已具备``痰瘀互结''的高危特征，应强制归入高风险；
    \item 若样本血脂正常（$N_{abn}=0$）但痰湿积分$\geq80$且活动能力评分$\text{ADL+IADL}\textless40$，则属``重度痰湿兼活动障碍''，同样强制归入高风险。
\end{enumerate}

记满足上述任一条件的人群为锚定高风险集，计算其综合风险评分均值$\bar{R}_{anchor}$。最终高风险阈值取统计分位数与临床锚点的较大值：
\begin{equation}
\theta_H = \max\bigl(Q_{0.75},\; \bar{R}_{anchor}\bigr)
\end{equation}
该设计确保：即使锚定人群的评分因数据分布原因低于75\%分位数，也不会被错误降层；同时避免了纯锚点法可能因锚定样本过少而导致阈值失真的问题。

\subsection{风险分层规则}
\begin{equation}
\text{风险等级}_i =
\begin{cases}
\text{低风险}, & R_i < \theta_L \\
\text{中风险}, & \theta_L \leq R_i < \theta_H \\
\text{高风险}, & R_i \geq \theta_H
\end{cases}
\end{equation}

\subsection{核心特征组合识别}
风险分层模型不仅能给出个体风险等级，还应能提炼出高风险人群的共性特征组合，为临床快速筛查提供规则化工具。本文采用\textbf{条件概率}（或称关联规则挖掘）方法，在全部特征空间中搜索满足以下条件的组合：
\begin{enumerate}
    \item 组合覆盖人数$|C|\geq 10$，避免过拟合到极端个案；
    \item 组合内高风险条件概率$P(\text{高风险}\mid C)\geq 50\%$，确保规则的临床可用性；
    \item 组合支持度（占总样本比例）$\geq 2\%$，保证规则具有一定普遍性。
\end{enumerate}

条件概率的计算公式为：
\begin{equation}
P(\text{高风险}\mid \text{特征组合 } C) = \frac{|\{i: R_i=\text{高}, i\in C\}|}{|C|}
\end{equation}

通过遍历候选特征（痰湿积分阈值、BMI阈值、TG阈值、活动能力阈值、血脂异常项数等）的笛卡尔积，筛选出满足上述三条准则的最强组合。表~\ref{tab:combo}展示了前三条核心规则。

\begin{table}[H]
\centering
\caption{痰湿体质高风险人群核心特征组合}
\label{tab:combo}
\begin{tabular}{clcc}
\toprule
\textbf{编号} & \textbf{特征组合$C$} & \textbf{高风险概率} & \textbf{覆盖人数} \\
\midrule
1 & 痰湿积分$\geq 60$ 且 活动总分$<40$ & 61.54\% & 26 \\
2 & 痰湿积分$\geq 60$ 且 BMI$\geq 24$ 且 TG$>1.7$ & 66.67\% & 18 \\
3 & 痰湿质为最高体质 且 血脂异常$\geq 2$项 & 53.76\% & 93 \\
\bottomrule
\end{tabular}
\end{table}

组合2的高风险概率最高（66.67\%），但覆盖人数较少，适用于精准筛查；组合3覆盖人数最多，适用于大规模初筛。临床实践中可先用组合3进行快速过滤，再用组合2对重点人群进行二次确认。

\subsection{模型验证：5折分层交叉验证与ROC分析}
为量化风险评分模型区分高风险与非高风险人群的能力，本文采用5折分层交叉验证（Stratified K-Fold CV）。分层策略确保每一折中高风险样本的比例与总体一致（约10.8\%），避免类别不平衡导致的评估偏差。

具体流程为：
\begin{enumerate}
    \item 将样本随机划分为5个子集；
    \item 依次取4个子集作为训练集，拟合Logistic回归模型；
    \item 在剩余1个子集（测试集）上预测发病概率，构造综合风险评分；
    \item 以测试集的真实高血脂标签（0/1）为金标准，计算准确率（Accuracy）、召回率（Recall）、F1分数及AUC。
\end{enumerate}

AUC（曲线下面积）衡量模型对正负样本的排序能力，取值0.5表示随机猜测，1.0表示完美区分。由于三级风险分层存在类别不平衡（高风险仅108例），传统准确率会受多数类主导而失真，因此AUC是更可靠的评价指标。

\subsection{权重敏感性分析}\label{sec:sens}
综合风险评分的权重$(w_1,w_2,w_3)=(0.5,0.3,0.2)$基于专家经验设定，其合理性需经敏感性检验。本文固定三级阈值$\theta_L,\theta_H$不变，将各权重独立扰动$\pm20\%$，观察高风险人群比例的变化幅度。若比例波动在$\pm5$个百分点以内，则认为权重选取稳健；反之则需引入熵权法等客观赋权手段。

\subsection{问题2结果与分析}

\subsubsection{风险分层结果}
经计算，样本综合风险评分$R_i$的30\%分位数为0.5749，75\%分位数为0.7750；锚定高风险集的评分为0.7213。取两者较大值，最终确定：
\begin{equation}
\theta_L=0.5749, \qquad \theta_H=0.7750
\end{equation}

三级风险分布如下：低风险300例（30.0\%）、中风险592例（59.2\%）、高风险108例（10.8\%）。该分布呈``橄榄型''，中等风险占据主体，符合流行病学中慢性病的普遍分布规律。

\subsubsection{高风险人群画像}
对108例高风险人群进行描述性统计，其典型特征为：
\begin{itemize}
    \item 痰湿积分均值57.0（标准差12.4），显著高于全样本均值（42.3）；
    \item TG均值2.56 mmol/L，约为全样本均值（1.89）的1.35倍；
    \item 血脂异常项数均值2.83，接近满项（4项）的70\%；
    \item 年龄分布以60$\sim$79岁组（年龄组3、4）为主，占比64.8\%。
\end{itemize}

\subsubsection{交叉验证与ROC结果}
5折交叉验证的AUC均值为$0.777$（各折在$0.738\sim0.797$之间），表明模型对高风险人群具有一定的区分能力，但尚未达到优秀水平（AUC$\textgreater0.85$）。这一结果提示：当前仅以5项血常规指标和痰湿积分为基础的风险评分，对中老年高血脂的预警能力仍有提升空间，后续可引入年龄、性别、生活方式等更多维度。

受类别不平衡影响，以0.5为默认阈值时的召回率极低（均值3.7\%），即模型倾向于将多数高风险样本判为非高风险。这在临床预警中不可接受——\textbf{宁可误报，不可漏报}。因此，实际应用中应将决策阈值下调（如降至0.25），以提升高风险召回率，同时结合临床锚定规则进行二次确认。

\subsubsection{权重敏感性分析结果}
表~\ref{tab:sens}展示了权重扰动对分层比例的影响。当$w_1$（血常规概率权重）从0.5升至0.6时，高风险比例从10.8\%跃升至24.3\%；当$w_1$降至0.4时，高风险比例跌至5.3\%。波动幅度超过$\pm10$个百分点，远超稳健性阈值（$\pm5$\%）。

\begin{table}[H]
\centering
\caption{风险评分权重敏感性分析}
\label{tab:sens}
\begin{tabular}{cccccc}
\toprule
\textbf{场景} & $w_1$ & $w_2$ & $w_3$ & \textbf{高风险比例} & \textbf{变化幅度} \\
\midrule
基准 & 0.50 & 0.30 & 0.20 & 10.8\% & --- \\
$w_1-20\%$ & 0.40 & 0.36 & 0.24 & 5.3\% & $-5.5$pp \\
$w_1+20\%$ & 0.60 & 0.24 & 0.16 & 24.3\% & $+13.5$pp \\
$w_2-20\%$ & 0.50 & 0.24 & 0.26 & 12.1\% & $+1.3$pp \\
$w_2+20\%$ & 0.50 & 0.36 & 0.14 & 10.7\% & $-0.1$pp \\
\bottomrule
\multicolumn{6}{l}{\small 注：pp = percentage points（百分点）。}
\end{tabular}
\end{table}

该结果表明：\textbf{风险评分对$w_1$（血常规概率权重）高度敏感}，而对$w_2,w_3$相对不敏感。其根源在于血常规概率$P_i$的样本方差远大于痰湿积分和血脂异常项数，权重的微小变化会被$P_i$的波动放大。改进方向为引入熵权法或CRITIC法客观赋权，降低主观经验的不确定性。

% ==========================================================
\section{问题3：干预方案优化模型}
% ==========================================================

\subsection{模型建立与分析}

问题3要求在中医调理原则、身体耐受度及经济成本的多重约束下，为痰湿体质患者制定6个月最优干预方案。这是一个典型的\textbf{带约束的组合优化问题}，其特殊性在于：部分决策变量（调理分级$x$）由患者初始痰湿积分强制确定，并非自由优化变量；真正的自由变量仅有活动强度$y$和每周频率$f$。

\subsubsection{决策变量与参数设定}
\begin{itemize}
    \item $x\in\{1,2,3\}$：中医调理分级。根据中医``辨证论治''原则，调理分级必须与痰湿积分区间严格匹配，不可跨级选择。
    \item $y\in\{1,2,3\}$：活动干预强度，对应轻度、中度、重度三档。
    \item $f\in\{1,2,\dots,10\}$：每周训练频率（次/周），受患者时间与体力限制，最多10次。
\end{itemize}

参数取值如下表所示：
\begin{table}[H]
\centering
\caption{干预成本参数表}
\label{tab:cost}
\begin{tabular}{cccc}
\toprule
\textbf{调理分级} $x$ & \textbf{月成本} $C_T(x)$（元/月） & \textbf{活动强度} $y$ & \textbf{单次成本} $C_A(y)$（元/次） \\
\midrule
1（基础） & 30 & 1（轻度） & 3 \\
2（中度） & 80 & 2（中度） & 5 \\
3（强化） & 130 & 3（重度） & 8 \\
\bottomrule
\end{tabular}
\end{table}

\subsubsection{约束条件}
\textbf{（1）调理分级约束（中医原则）}

中医强调``同病异治、异病同治''，调理方案必须与体质严重程度匹配：
\begin{equation}
x =
\begin{cases}
1, & S_0 \leq 58 \\
2, & 59 \leq S_0 \leq 61 \\
3, & S_0 \geq 62
\end{cases}
\end{equation}
该约束为硬约束（Hard Constraint），在优化过程中不可违反。

\textbf{（2）年龄耐受约束}

中老年人群各年龄段的心肺功能、骨骼肌肉状态差异显著，活动强度必须受年龄限制：
\begin{equation}
y \leq y_{max}^{age} =
\begin{cases}
3, & A\in\{1,2\}\ (40\sim59\text{岁}) \\
2, & A\in\{3,4\}\ (60\sim79\text{岁}) \\
1, & A=5\ (80\sim89\text{岁})
\end{cases}
\end{equation}

\textbf{（3）活动能力耐受约束}

ADL（日常生活活动）与IADL（工具性日常生活活动）量表总分$M$直接反映患者的运动耐受能力。即使年龄较轻，若活动能力评分过低，也应降低运动强度以防跌倒或心血管意外：
\begin{equation}
y \leq y_{max}^{mob} =
\begin{cases}
1, & M < 40 \\
2, & 40 \leq M < 60 \\
3, & M \geq 60
\end{cases}
\end{equation}
综合强度上限取两者最小值：$y \leq \min(y_{max}^{age}, y_{max}^{mob})$。该设计体现了``木桶效应''——患者的活动能力上限由其最弱的生理维度决定。

\textbf{（4）频率约束与成本约束}

每周训练频率必须在合理范围内：$1 \leq f \leq 10$。

总成本包括6个月的调理费用和24周（每月按4周计）的活动费用：
\begin{equation}
\text{Cost} = 6\cdot C_T(x) + 24\cdot f\cdot C_A(y) \leq 2000\quad(\text{元})
\end{equation}

\subsubsection{目标函数与依从性约束}
优化的核心目标是最大程度降低痰湿积分。根据题目假设，每月痰湿积分下降率$r$由活动强度$y$和频率$f$共同决定。但题目给出的线性公式未考虑老年医学中的\textbf{干预疲劳效应}与\textbf{依从性衰减}——对80--89岁或活动能力低下的患者，每周高频干预不仅难以坚持，还可能因过度疲劳导致效果打折甚至引发安全问题。

因此，本文在原始公式基础上引入依从性衰减函数$\eta(f; A, M)$：
\begin{equation}
r =
\begin{cases}
0, & f < 5 \\
0.03\cdot(y-1) + 0.01\cdot\bigl(\eta(f;A,M)-5\bigr), & f \geq 5
\end{cases}
\end{equation}
其中有效频率$\eta(f;A,M)$为：
\begin{equation}
\eta(f;A,M) =
\begin{cases}
\min\bigl\{5+0.5(f-5),\;6.0\bigr\}, & A=5\ (80\sim89\text{岁}) \\
\min\bigl\{5+0.6(f-5),\;6.2\bigr\}, & M<40 \quad \text{（活动能力低下）} \\
5+0.8(f-5), & A\in\{3,4\}\ (60\sim79\text{岁}) \\
f, & A\in\{1,2\}\ \text{且}\ M\geq40
\end{cases}
\end{equation}

该设计的临床依据为：
\begin{itemize}
    \item \textbf{频率门槛效应}：每周训练少于5次时，积分基本不降，提示运动干预必须达到一定剂量才起效；
    \item \textbf{依从性衰减}：高龄或低活动能力患者超出5次/周的部分，有效频率按比例打折（50\%--80\%），体现生理耐受极限；
    \item \textbf{疲劳封顶}：80--89岁及活动能力低下者的有效频率分别封顶6.0和6.2次/周，即使名义频率$f=10$，实际等效仅约6次，防止模型推荐不切实际的极端方案。
\end{itemize}

6个月后痰湿积分为：
\begin{equation}
S_6 = S_0 \cdot (1-r)^6
\end{equation}
取对数后可见，$\ln S_6$与$r$呈线性关系，最小化$S_6$等价于最大化$r$。当多个方案$S_6$相同时（通常发生在$r$触及封顶后），优先选择总成本较低者，以节约医疗资源。

\subsection{求解方法：穷举法}
对单名患者而言，调理分级$x$已由$S_0$强制确定，活动强度$y$的上界由年龄和活动能力共同约束（最多3级），频率$f$取1$\sim$10的整数值。因此，单名患者的合法组合数极为有限：
\begin{equation}
N_{comb} = y_{max} \times 10 \leq 3 \times 10 = 30
\end{equation}

在如此小的决策空间中，任何启发式算法（如遗传算法、模拟退火）均无必要，直接采用\textbf{穷举法（Brute-force Enumeration）}即可在毫秒级求得全局最优解。算法流程如下：

\begin{enumerate}
    \item \textbf{输入}：患者初始痰湿积分$S_0$、年龄组$A$、活动总分$M$；
    \item 根据$S_0$确定调理分级$x$，根据$A$和$M$确定$y_{max}$；
    \item \textbf{枚举}所有满足$1\leq y\leq y_{max}$、$1\leq f\leq 10$的$(y,f)$组合；
    \item 对每个组合计算月下降率$r$、6个月后积分$S_6$、总成本Cost；
    \item 剔除Cost$\textgreater2000$的非法组合；
    \item 在合法组合中选取$S_6$最小者；若$S_6$相同，取Cost最小者；
    \item \textbf{输出}：最优方案$(x^*,y^*,f^*)$及对应的$S_6^*$、Cost$^*$。
\end{enumerate}

穷举法的优势在于：①保证全局最优，不存在陷入局部极值的风险；②实现简单，代码可读性强；③便于敏感性分析——通过穷举结果可直接绘制$S_6$随$(y,f)$变化的响应曲面。

\subsection{问题3结果}

\subsubsection{样本ID 1、2、3最优方案}

\begin{table}[H]
\centering
\caption{样本ID 1、2、3最优干预方案}
\resizebox{\textwidth}{!}{%
\begin{tabular}{ccccccccc}
\toprule
\textbf{样本ID} & \textbf{痰湿积分} & \textbf{年龄组} & \textbf{活动总分} & \textbf{调理分级} & \textbf{活动强度} & \textbf{每周频率} & \textbf{总成本} & \textbf{6个月后积分} \\
\midrule
1 & 64 & 2(50$\sim$59) & 38 & 3级(强化) & 1级 & 7次 & 1284元 & 59.53 \\
2 & 58 & 1(40$\sim$49) & 40 & 1级(基础) & 2级 & 10次 & 1380元 & 35.17 \\
3 & 59 & 1(40$\sim$49) & 63 & 2级(中度) & 2级 & 10次 & 1680元 & 35.77 \\
\bottomrule
\end{tabular}%
}
\end{table}

\subsubsection{``患者特征—最优方案''匹配规律}
为揭示不同患者特征与最优方案之间的系统性对应关系，对全部痰湿体质患者（$n=312$）执行穷举优化，并按活动能力、年龄组、痰湿积分进行分组统计，归纳出以下规律：

\textbf{（1）活动能力决定强度天花板，频率受依从性约束}

在引入依从性衰减后，最优频率不再无条件拉满至10次/周，而是受年龄和活动能力的双重约束：
\begin{itemize}
    \item \textbf{活动能力差}（$M<40$）：强度被锁死在1级；同时因依从性衰减，有效频率封顶约6.2次/周，最优名义频率$f^*\approx7.2$次/周，每月下降率$r\approx1.2\%$，6个月后积分约为$S_0\times0.988^6\approx0.93\,S_0$。
    \item \textbf{活动能力中等}（$40\leq M<60$）：可选2级强度，无频率封顶（40--59岁）或80\%衰减（60--79岁），多数患者$f^*\approx9.5$次/周，$r\approx8\%$，6个月后积分约为$0.92^6\approx0.606\,S_0$。
    \item \textbf{活动能力好}（$M\geq 60$）：$y_{max}^{mob}=3$，但多数情况下$y^*=2$而非3。原因在于：从2级升至3级，单次活动成本从5元增至8元，$r$仅从8\%增至11\%，积分下降收益有限，而成本显著上升（24周$\times10$次$\times3$元差额$=720$元），性价比不如维持2级。
\end{itemize}

\textbf{（2）高龄患者以``稳''为主}

年龄组5（80$\sim$89岁）的患者，$y_{max}^{age}=1$，且依从性衰减最强（超出5次部分仅50\%有效，封顶6.0次）。因此，这类患者的最优名义频率$f^*$约为7次/周，实际等效频率仅约6次/周，成本约为$6\times C_T(x)+24\times7\times3=6C_T(x)+504$元，通常在1000$\sim$1300元之间。这体现了老年医学``低强度、高安全''的干预原则。

\textbf{（3）调理分级对成本的刚性影响}

由于调理分级$x$由痰湿积分强制确定，高积分患者必须承担更高的固定月费。以$f=10$、$y=2$为例：
\begin{itemize}
    \item $x=1$（基础）：月费30元，6个月调理费180元；
    \item $x=2$（中度）：月费80元，6个月调理费480元；
    \item $x=3$（强化）：月费130元，6个月调理费780元。
\end{itemize}
因此，在相同活动方案下，痰湿积分从58分升至62分（跨越阈值），总成本将跃增300元。这一``阶梯式''成本结构提示：在积分临界点附近，通过预干预降低痰湿积分以实现调理分级``降档''，具有显著的经济学价值。

\textbf{（4）成本—效果分析}

对全部痰湿体质患者（$n=312$）执行优化后，按实际调理分级统计平均降分幅度与平均成本，计算单位成本降分效果：
\begin{itemize}
    \item $x=1$（基础调理）：平均降分16.54分，平均成本1215.12元，单位成本降分0.0136分/元；
    \item $x=2$（中度调理）：平均降分15.85分，平均成本1442.73元，单位成本降分0.0110分/元；
    \item $x=3$（强化调理）：平均降分16.35分，平均成本1724.80元，单位成本降分0.0095分/元。
\end{itemize}

可见，基础调理（$x=1$）的单位成本降分效果最优，这与中医``轻证轻治''的原则相吻合。需要说明的是，强化调理（$x=3$）的患者群体本身痰湿积分更高（均值$\geq62$），依从性约束更强（活动能力差者占比更高），导致其平均降分幅度并未显著高于中度调理，而固定月费更高，拉低了单位成本效果。这提示：在痰湿积分临界点附近（58--62分），通过预干预实现调理分级``降档''，可在不牺牲疗效的前提下显著降低医疗支出。

% ==========================================================
\section{模型评价与改进}
% ==========================================================

\subsection{模型优点}
\begin{enumerate}
    \item \textbf{可解释性强}：采用Logistic回归、交叉验证、穷举法等传统统计方法，结果具有明确的临床解释性。
    \item \textbf{计算高效}：问题3决策空间仅30种组合，穷举法毫秒级求解，无需启发式算法。
    \item \textbf{阈值选取兼顾统计与临床}：风险分层既参考了数据分布，又锚定了临床先验，避免纯数据驱动的极端化。
\end{enumerate}

\subsection{模型缺点}
\begin{enumerate}
    \item \textbf{风险评分权重主观性高}：本文采用专家经验权重$(0.5,0.3,0.2)$，敏感性分析显示高风险比例对$w_1$扰动高度敏感（波动$\pm13.5$pp），主观赋权存在较大不确定性。
    \item \textbf{线性假设局限}：逐步回归与Logistic模型假设变量间为线性关系，可能遗漏非线性交互效应（如痰湿积分$\times$TG的交互项）。
    \item \textbf{干预效果参数未标定}：问题3中下降率公式（3\%/级、1\%/次）为题目给定假设，未基于随访数据标定，实际个体间差异可能显著。
    \item \textbf{交叉验证召回率偏低}：默认0.5阈值下高风险召回率仅3.7\%，尽管AUC达0.777，但临床预警要求``宁可误报、不可漏报''，当前阈值策略有待优化。
\end{enumerate}

\subsection{模型改进方向}
\begin{enumerate}
    \item \textbf{客观赋权}：引入CRITIC熵权法或基于变异系数的客观赋权，替代主观经验权重，降低分层结果对权重选取的敏感度。
    \item \textbf{非线性建模}：采用广义加性模型(GAM)或梯度提升树(GBDT)捕捉痰湿积分与血脂指标间的非线性交互效应。
    \item \textbf{概率性动态干预模型}：将确定性下降率改为Logistic饱和曲线$r(f)=r_{\max}/(1+e^{-k(f-f_0)})$，反映运动干预的边际递减效应；同时引入依从性概率$\pi(f)=\max(0,1-\lambda f)$，使模型更贴近临床实际。
    \item \textbf{阈值优化}：以最大化F$\beta$分数（$\beta=2$，侧重召回率）为目标，在训练集上搜索最优决策阈值，而非固定0.5或分位数。
\end{enumerate}

% ==========================================================

\newpage


% ==========================================================

\begin{thebibliography}{99}
\bibitem{dong2004algo}
董乘宇. 数学建模竞赛中应当掌握的十类算法[J]. 数模, 2004, 1(1): 12-14.
\bibitem{wang2009tcm}
王琦. 中医体质学[M]. 北京: 人民卫生出版社, 2009.
\bibitem{zhinan2016}
中国成人血脂异常防治指南修订联合委员会. 中国成人血脂异常防治指南(2016年修订版)[J]. 中华心血管病杂志, 2016, 44(10): 833-853.
\end{thebibliography}

% ==========================================================

\newpage
\appendix
\section{附录}
% ==========================================================

\subsection{软件及运行环境说明}

本论文所有数据分析、模型求解及可视化均基于以下软件环境完成：

\begin{table}[H]
\centering
\caption{软件及环境列表}
\begin{tabular}{lll}
\toprule
\textbf{类别} & \textbf{名称} & \textbf{版本/说明} \\
\midrule
操作系统 & macOS / Windows 10 & 64位 \\
编程语言 & Python & 3.12+ \\
数据分析库 & pandas, numpy & 2.x, 1.26+ \\
机器学习库 & scikit-learn & 1.4+ \\
统计建模库 & statsmodels, scipy & 0.14+, 1.12+ \\
文档排版 & \LaTeX & TeX Live 2024 \\
\bottomrule
\end{tabular}
\end{table}

\subsection{主要计算机源程序}

本论文的全部源代码使用 Python 3.12 编写，保存为文件 \texttt{solution.py}，随支撑材料一并提交。程序涵盖数据读取、预处理、问题1（相关性分析+组间差异检验+逐步Logistic回归+九种体质贡献度分析+VIF计算）、问题2（风险评分构造+阈值划分+5折交叉验证+ROC分析+权重敏感性分析）、问题3（带依从性约束的穷举优化+结果汇总+可视化）的完整流程。下面给出问题3中\textbf{带依从性约束的穷举优化算法}的核心代码片段，其余模块详见支撑材料。

\begin{lstlisting}[language=Python,caption={问题3 带依从性约束的穷举优化核心代码}]
def calc_monthly_drop_rate(y, f, age_group, mobility):
    """计算每月痰湿积分下降率（含依从性衰减）"""
    if f < 5:
        return 0.0
    # 依从性衰减 + 疲劳封顶
    if age_group == 5:          # 80-89岁
        effective_f = min(5 + (f - 5) * 0.5, 6.0)
    elif mobility < 40:         # 活动能力低下
        effective_f = min(5 + (f - 5) * 0.6, 6.2)
    elif age_group in [3, 4]:   # 60-79岁
        effective_f = 5 + (f - 5) * 0.8
    else:                       # 40-59岁且活动能力好
        effective_f = f
    base_drop = 0.03 * (y - 1)
    extra_drop = 0.01 * (effective_f - 5)
    return base_drop + extra_drop

def optimize_intervention(patient):
    """对单名患者穷举求解最优干预方案"""
    s0 = patient['phlegm']
    x  = get_treatment_level(s0)        # 由痰湿积分强制匹配
    y_max = get_max_activity_intensity(
                patient['age_group'],
                patient['mobility'])    # 年龄与活动能力双约束

    best_plan, best_final = None, float('inf')
    for y in range(1, y_max + 1):       # 枚举活动强度
        for f in range(1, 11):          # 枚举每周频率
            r = calc_monthly_drop_rate(
                    y, f,
                    patient['age_group'],
                    patient['mobility'])
            s_final = s0 * ((1 - r) ** 6)
            total_cost = 6 * C_TREAT[x] + 24 * f * C_ACT[y]
            if total_cost > 2000:
                continue                # 成本硬约束
            if s_final < best_final - 1e-6:
                best_final = s_final
                best_plan = (x, y, f, total_cost, s_final)
            elif abs(s_final - best_final) < 1e-6 \
                    and total_cost < best_plan[3]:
                best_plan = (x, y, f, total_cost, s_final)
    return best_plan
\end{lstlisting}

\subsection{AI使用情况说明}

在本竞赛论文的撰写过程中，使用了生成式人工智能（AI）工具辅助完成以下工作：

\begin{enumerate}
    \item \textbf{文献综述与知识检索}：利用AI快速检索中医体质学、Logistic回归及Fisher判别分析的相关定义与公式表述，确保学术用语准确。
    \item \textbf{代码框架生成}：借助AI生成Python代码的初始框架（数据读取、函数模板），随后由参赛队员根据题目约束条件自行修改参数、补全业务逻辑并调试运行。
    \item \textbf{排版与格式检查}：使用AI辅助将文本内容转换为 \LaTeX 格式，并依据竞赛排版规范进行页码、目录、参考文献等结构调整。
\end{enumerate}

所有核心建模思路、数学推导、结果分析以及代码的最终调试均由参赛队员独立完成，AI仅作为效率辅助工具。参赛者已对所有内容的真实性、准确性和原创性负责。

\end{document}
